\section{Objekterkennungsmodell}


\subsection{Media Pipe Model Maker}
Eine andere Möglichkeit den Tensorflow Model Maker zu umgehen ist der Verwendung des Media Pipe Model Makers von Google. Durchgeführt mit Google Colab, da hier eine kostenlose GPU zur Verfügungsteht und das Training damit deutlich schneller durchgeführt werden kann.
\begin{lstlisting}
    !python --version
    !pip install --upgrade pip
    !pip install mediapipe-model-maker
\end{lstlisting}

\begin{lstlisting}
    from google.colab import files
    import os
    import json
    import tensorflow as tf
    assert tf.__version__.startswith('2')

    from mediapipe_model_maker import object_detector
\end{lstlisting}

Es muss der Datensatz an Bilder im Colab hochgeladen sein. In diesem Beispiel liegt ein Datenset im Pascal VOC-Format vor und heißt \texttt{main\_dataset.zip}. Hierbei ist es auch wichtig auf die richtige Ordnerstruktur zu achten, damit der Model Maker die Bilder und Labels korrekt erkennt. Die Struktur ist wie folgt:
\begin{lstlisting}
    main_dataset/
        train/
            images/
                image1.jpg
                image2.jpg
            Annotations/
                image1.xml
                image2.xml
        valid/
            images/
                image3.jpg
                image4.jpg
            Annotations/
                image3.xml
                image4.xml
\end{lstlisting}


\begin{lstlisting}
    !unzip main_dataset.zip

    train_dataset_path = "main_dataset/train"
    validation_dataset_path = "main_dataset/valid"
\end{lstlisting}

\begin{lstlisting}
     train_data = object_detector.Dataset.from_pascal_voc_folder(
        'main_dataset/train',
        cache_dir="/tmp/od_data/train",
    )

    validation_data = object_detector.Dataset.from_pascal_voc_folder(
        'main_dataset/valid',
        cache_dir="/tmp/od_data/validation")

\end{lstlisting}

\begin{lstlisting}
    hparams = object_detector.HParams(batch_size=8, learning_rate=0.3, epochs=50, export_dir='exported_model')

    options = object_detector.ObjectDetectorOptions(
        supported_model=object_detector.SupportedModels.MOBILENET_V2,
        hparams=hparams
    )

    model = object_detector.ObjectDetector.create(
        train_data=train_data,
        validation_data=validation_data,
        options=options
    )

\end{lstlisting}

\begin{lstlisting}
    train_data = object_detector.Dataset.from_pascal_voc_folder(train_dataset_path, cache_dir="/tmp/od_data/train")
    validation_data = object_detector.Dataset.from_pascal_voc_folder(validation_dataset_path, cache_dir="/tmp/od_data/validation")
    print("train_data size: ", train_data.size)
    print("validation_data size: ", validation_data.size)
\end{lstlisting}

Anschließend kann das Model exportiert und heruntergeladen werden. Dabei wird das Modell in das \texttt{.tflite}-Format konvertiert, welches für die Android App benötigt wird. Der Export erfolgt mit dem folgenden Code:
\begin{lstlisting}
    model.export_model()
    !ls exported_model
    files.download('exported_model/model.tflite')
\end{lstlisting}